% SPDX-License-Identifier: CC-BY-4.0
% Copyright (c) 2026 Martin Schröder <info@swedishembedded.com>

\section{Results}\label{sec:results}

Every exam in this section is \textbf{pilot scale}: 31 to 42 tasks over 6 or 7 families. Their families are fixed by each run's held-out split before the candidate is trained,
but they are far too small to carry a claim, and
several candidates were selected against the same task sets (two exams share a 35-task set,
three share a 31-task set), so the comparisons are not independent. None is a frozen final test.
\Cref{sec:frozen} describes the powered exam that is.

\subsection{Candidates}

\Cref{tab:training} lists the training runs. The runs differ in more than the data (the
recipe changed between them, \cref{sec:training-defects}); we do not claim that any difference in
outcome is due to a single change. Held-out loss and token accuracy fall in every run, but token
accuracy is nearly flat (for v7 it moves from 0.731 to 0.750), which is expected for a stylistic
change~\citep{lin2024urial} and carries no information about voice. Losses are
not comparable across runs with different held-out sets.

\begin{table}[t]
\centering\small
\resizebox{\textwidth}{!}{\input{generated/training_table}}
\caption{Training runs. ``Step kept'' is the step of the adapter carried (``last'' for
unmonitored runs that carry the final step). v6 trained on dialogues only; ablation and v7 on
dialogues plus 210 fixed-prompt voice rows; e4 on dialogues plus template-prompted voice
chunks under the current recipe; e1 on dialogues plus the 210 fixed-prompt voice rows under the same recipe (not yet examined); adams on Adams dialogues plus voice rows. Peak learning rate
was not recorded for the earlier runs. Held-out records are scored with the base before
and the adapter after.}
\label{tab:training}
\end{table}

\begin{figure}[t]
\centering
\begin{tikzpicture}
\begin{groupplot}[group style={group size=2 by 1, horizontal sep=1.6cm}, width=0.48\textwidth,
  height=5.2cm, tick label style={font=\scriptsize}, label style={font=\small}, grid=major,
  grid style={gray!25}, legend style={font=\scriptsize}]
\nextgroupplot[xlabel={step}, ylabel={learning rate ($10^{-4}$)}, ymin=0, ymax=3.3, legend pos=south west]
\addplot[thick, red!70!black, mark=*, mark size=1.2pt] table[x=step, y expr=\thisrow{lr}*10000, col sep=comma] {generated/curve_v7.csv};
\addlegendentry{v7}
\addplot[thick, blue!60!black, mark=square*, mark size=1.2pt] table[x=step, y expr=\thisrow{lr}*10000, col sep=comma] {generated/curve_e4.csv};
\addlegendentry{e4}
\nextgroupplot[xlabel={step}, ylabel={monitoring loss}, ymin=0.85, ymax=1.3, legend pos=north east]
\addplot[thick, red!70!black, mark=*, mark size=1.2pt] table[x=step, y=monitor, col sep=comma] {generated/curve_v7.csv};
\addlegendentry{v7}
\addplot[thick, blue!60!black, mark=square*, mark size=1.2pt] table[x=step, y=monitor, col sep=comma] {generated/curve_e4.csv};
\addlegendentry{e4}
\end{groupplot}
\end{tikzpicture}
\caption{Learning rate and monitoring loss at each evaluation. v7 (recipe before the
corrections) stopped at step 56 of 105 with the rate still at $2.05\times10^{-4}$ and kept step
28; e4 (current recipe) annealed to $2\times10^{-5}$ and kept step 105. The monitoring sets
differ (v7: one family, dialogue records only; e4: several families including the writer's
text), so the levels are not comparable, only the shapes.}
\label{fig:curves}
\end{figure}

\subsection{Pilot-scale exams}

\Cref{tab:exams} gives, for each exam, the number of tasks judged right in each arm, the number
of invented-specifics answers, and the family-level sign test of the deployed candidate
against the unprompted and the prompted base. The $p$-values are one-sided exact tests
over discordant families, as the exam logged them.

\begin{table}[t]
\centering\scriptsize
\setlength{\tabcolsep}{3pt}
\resizebox{\textwidth}{!}{\input{generated/exams_table}}
\caption{Pilot-scale exams. ``Tasks/fam.'' is tasks over families; ``prompt.'' is the base under
the persona prompt; ``cand.'' is the candidate under the persona prompt; ``invented'' counts
answers that state a number or name absent from the source (of all tasks in the column).
Families are counted once whatever their size. v6 and ablation share one task set and v7,
rehearsal and e4 share another, so the pooled row, which sums the four rows v6, ablation,
v7 and adams, counts the first task set twice and is descriptive only.}
\label{tab:exams}
\end{table}

\begin{figure}[t]
\centering
\begin{tikzpicture}
\begin{axis}[ybar, width=\textwidth, height=5.8cm, bar width=7pt, ymin=0, ymax=0.7,
  ylabel={share of tasks judged right}, symbolic x coords={v6,ablation,v7,rehearsal,e4,adams},
  xtick=data, tick label style={font=\scriptsize}, label style={font=\small}, enlarge x limits=0.1,
  legend style={font=\scriptsize}, legend pos=north west, ymajorgrids=true,
  grid style={gray!25}]
\addplot[fill=gray!55] table[x=label, y=base, col sep=comma] {generated/arms.csv};
\addplot[fill=orange!70] table[x=label, y=prompted, col sep=comma] {generated/arms.csv};
\addplot[fill=blue!60!black] table[x=label, y=candidate, col sep=comma] {generated/arms.csv};
\legend{base, prompted base, candidate (persona prompt)}
\end{axis}
\end{tikzpicture}
\caption{Share of tasks judged right by arm in each pilot-scale exam (31 to 42 tasks each).
The prompt closes more of the gap than the fine-tune does; no panel has the power to resolve the
remaining difference.}
\label{fig:arms}
\end{figure}

The pattern is consistent but not conclusive.
\begin{itemize}
\item \textbf{The persona prompt does most of the work.} Pooled over the four exams
  (\PooledTasks\ tasks), the unprompted base is judged right on \PooledBase\ tasks
  (\PooledBasePct\%) and the prompted base on \PooledPrompted\ (\PooledPromptedPct\%): ten points
  from a prompt. The adapter is judged right on \PooledCand\ (\PooledCandPct\%), nine points
  above the prompted base.
\item \textbf{Invented specifics fall with the prompt and, a little, again with the
  adapter.} Pooled, the unprompted base invents on \PooledBaseInv\ answers, the prompted
  base on \PooledPromptedInv\ and the adapter on \PooledCandInv.
\item \textbf{No candidate separates from the prompted base.} The family-level test against
  the prompted base has $p\geq0.109$ in every exam. The best case, the ablation candidate
  (voice rows, 18 against 14 right), has five wins in six discordant families. The task-level
  bound (\cref{tab:bounds}) shows that the same exam could not have reached significance at the
  task level under any split of its discordant tasks, since only four net tasks separate the two
  arms.
\item \textbf{Against the unprompted base the evidence is stronger but still thin.} The
  ablation candidate won all five of its discordant families ($p=0.031$, one-sided; the test
  is on five units), the only family-level result below 0.05 in the table. Because the
  comparison is with an unprompted base, it does not separate what the prompt and what the
  adapter contribute.
\item \textbf{Adams is the largest difference.} 17 against 9 right with the prompted base
  (net $+8$ tasks) and 6 against 8 invented. At the task level the one-sided $p$ could lie
  anywhere from 0.004 to 0.084; at the family level, with six families and three discordant,
  the logged value is 0.125.
\item \textbf{Rehearsal and the new recipe do not change the picture.} The rehearsal candidate
  (v7 data plus replay of the base's own answers) judged 14 of 31 right against 13 for the
  prompted base; the e4 candidate 16 against 13, with 3 invented against 6 for the prompted
  base. Neither is distinguishable (family-level $p=0.5$ and $0.125$).
\end{itemize}

\begin{table}[t]
\centering\small
\input{generated/bounds_table}
\caption{What each exam's marginals allow: with $m$ net tasks in the candidate's favour over the
prompted base, the smallest one-sided exact task-level sign $p$ over every possible split of
discordant tasks is $0.5^m$ (ignoring family clustering, which only raises it).}
\label{tab:bounds}
\end{table}

\paragraph{The honest reading.}
The exams are consistent with an adapter that is a few points better than the prompted base
and with one that is no better. Pooled counts of tasks are descriptive, not a test: the
pooled task sets overlap, families are clustered, and the exams were used repeatedly to choose
recipes (\cref{sec:threats}). We therefore do \emph{not} claim that fine-tuning beats prompting
for this task. What the data do support is narrower: a prompt alone changes the answers
a good deal; a fine-tune that sees the author's text makes fewer unsupported specific claims
than the unprompted base (\PooledCandInv\ against \PooledBaseInv\ of \PooledTasks\ answers), and
is no worse than the prompted base on the measures we have.

\subsection{Release gate}

No Qwen3 candidate has been released. The Jefferson v7 candidate passed the improvement check
(subject to the defect above: 33 wins against 3, with 12 of 12 discordant families) and failed
the anchor check (accuracy 0.773 against 0.800, drop 0.027 against a bound of 0.02,
\cref{tab:anchor}); the serve check was not measured (\cref{sec:limits}). The Adams candidate
passed improvement (54 wins against 9 over 207 tasks, 165 of them variants; 13 of 13 discordant
families) and failed the anchor (0.764 against 0.800). The rehearsal candidate
passed the anchor (0.818 against 0.800, 9 candidate-only against 7 base-only items) but failed
improvement (29 wins against 7 items, but 8 of 11 discordant families, $p=0.113$) and the serve
check, which at that time counted ungraded served answers as different (\cref{sec:residency}).
A release was made once, for an earlier student model (a DeepSeek-R1 distilled 7B), where the
served model re-answered 8 of 8 sampled tasks alike. The rehearsal result is one run of one seed: it
does not show that rehearsal protects the anchor, because the anchor's own resolution is
only $\pm0.037$ (\cref{sec:power}).

\subsection{The frozen exam}\label{sec:frozen}

The final test and dev suite were reserved from the pool of letters left out of the materials,
before any task was generated. Fifty families were reserved for each. Generation admitted tasks
from only 39 families of the final test (244 tasks: 22 families with the full eight, the
others with one to seven) and from 44 of the dev suite (166 tasks), in the three kinds
conversation, advice and explanation. The planned 50$\times$8 design (400 tasks) was therefore
not reached. At the realised size the simulated power for a ten-point gain is 0.53 under the
planning assumptions (\cref{tab:power}). A 25-family pilot (155 tasks) is the first run on it; it
estimates the discordance and intraclass correlation that the simulation currently assumes.
\Cref{tab:pilot} lists the pilots. Runs that had not completed when this paper was written are
marked; they are not zeros, they were not measured.

\begin{table}[t]
\centering\small
\resizebox{\textwidth}{!}{\input{generated/pilot_table}}
\caption{Pilots on the frozen exam (deployed candidate against the prompted arm of the v7
pilot, paired by task, difference in points with the family-clustered 95\% interval). To update
a row, edit \code{data/pilot.csv} and rebuild.}
\label{tab:pilot}
\end{table}

\subsection{Cost}

Reported for planning, not as performance claims. Generating tasks over 19 of 1,629 parts took
6,477~s; teaching, verification, authoring, variants and selection for the v7 run took about five hours together
(257 teacher solves, of which 161 verified and 95 failed); training v7 took about two hours for 56
steps; an exam of 31 tasks over three arms with judging took about 38 minutes; a release gate
over 259 tasks and two arms took about two hours and 12 minutes. The first two figures are
the reason the pipeline cannot cover a corpus of this size with generated tasks (covering all
1,629 parts at the observed pace is 555,322~s) and the reason the voice view, which needs no
model, carries most of the author's text.
